% ============================================================
% BÖLÜM 2: PROBLEM DURUMU
% Genişletilmiş ve Akademik Referanslarla Zenginleştirilmiş
% ============================================================

\chapter{PROBLEM DURUMU}

\section{Mevcut Durumun Betimlenmesi}

Türkiye'de okul psikolojik danışmanlığı hizmetleri, Milli Eğitim Bakanlığı bünyesinde yürütülmektedir. Rehberlik ve Araştırma Merkezleri (RAM) ile okullarda görevli psikolojik danışmanlar, öğrencilerin eğitsel, mesleki ve kişisel-sosyal gelişimlerini desteklemekle yükümlüdür. Ancak mevcut sistemde çeşitli yapısal sorunlar bulunmaktadır.

\subsection{Danışman-Öğrenci Oranı Sorunu}

Uluslararası standartlara göre ideal danışman-öğrenci oranı 1:250 olarak kabul edilmektedir (American School Counselor Association, 2019). UNESCO (2021) raporuna göre, gelişmiş ülkelerde bu oran 1:300 ile 1:500 arasında değişmektedir. Türkiye'de ise bu oran çoğu zaman 1:1000'in üzerindedir. MEB (2020) verilerine göre, bazı bölgelerde bu oran 1:1500'e kadar çıkabilmektedir. Bu durum, bireysel değerlendirme ve müdahale süreçlerini ciddi biçimde kısıtlamaktadır.

\begin{table}[H]
\centering
\caption{Uluslararası Karşılaştırma: Danışman-Öğrenci Oranları}
\label{tab:danismanoran}
\begin{tabular}{lcc}
\toprule
\textbf{Ülke/Kuruluş} & \textbf{Danışman-Öğrenci Oranı} & \textbf{Kaynak} \\
\midrule
Amerika Birleşik Devletleri & 1:424 & ASCA, 2021 \\
İngiltere & 1:350 & DfE, 2020 \\
Almanya & 1:300 & KMK, 2019 \\
Finlandiya & 1:200 & OPH, 2020 \\
Japonya & 1:400 & MEXT, 2021 \\
Türkiye (Ortalama) & 1:800+ & MEB, 2020 \\
ASCA Önerisi & 1:250 & ASCA, 2019 \\
\bottomrule
\end{tabular}
\end{table}

Bu veriler, Türkiye'deki okul psikolojik danışmanlarının uluslararası standartların çok üzerinde bir iş yükü ile karşı karşıya olduğunu ortaya koymaktadır.

\subsection{Projektif Testlerin Kullanım Zorlukları}

Projektif çizim testleri, çocuk değerlendirmesinde yaygın biçimde kullanılmakta; ancak çeşitli zorluklarla karşılaşılmaktadır:

\begin{enumerate}
  \item \textbf{Yorumlama Karmaşıklığı:} HTP testinin yorumlanması, çok sayıda sembolik anlamın değerlendirilmesini gerektirmektedir. Buck (1948) ve Hammer (1958), yalnızca ``ev'' çiziminde kapı, pencere, çatı, baca, duvar kalınlığı, konumlandırma gibi onlarca öğenin ayrı ayrı değerlendirilmesi gerektiğini belirtmişlerdir. Jolles (1971), HTP yorumlama kataloğunda 400'ü aşkın sembolik gösterge tanımlamıştır.

  \item \textbf{Eğitim Yetersizliği:} Lisans düzeyinde PDR eğitimi, projektif testler konusunda sınırlı kalmakta; derinlemesine uzmanlık için lisansüstü eğitim veya özel sertifika programları gerekmektedir (Özgüven, 2011). YÖK (2018) müfredatlarının incelenmesinde, projektif test eğitiminin genellikle 1-2 dersten ibaret olduğu görülmektedir.

  \item \textbf{Öznel Değerlendirme:} Aynı çizim, farklı değerlendiriciler tarafından farklı biçimlerde yorumlanabilmektedir. Bu durum, değerlendirmeler arası güvenilirlik sorunlarına yol açmaktadır. Handler ve Habenicht (1994), HTP testinin değerlendiriciler arası güvenilirliğinin .31 ile .96 arasında değiştiğini raporlamıştır.

  \item \textbf{Zaman Gereksinimleri:} Tam bir HTP değerlendirmesi, deneyimli bir klinisyen için bile 45-60 dakika sürebilmektedir (Groth-Marnat, 2009). Okul ortamında bu süre çoğu zaman ayrılamamaktadır.
\end{enumerate}

\subsection{Teknoloji Entegrasyonu Eksikliği}

Eğitim alanında teknoloji kullanımı hızla yaygınlaşırken, psikolojik danışmanlık hizmetlerinde dijital araçların kullanımı görece sınırlı kalmıştır. COVID-19 pandemisi, çevrimiçi danışmanlık hizmetlerinin yaygınlaşmasını hızlandırmış olsa da (Ribeiro ve ark., 2021), özellikle değerlendirme süreçlerinde yapay zekâ destekli araçların kullanımı, PDR alanında henüz yaygınlaşmamıştır.

Dünya Sağlık Örgütü (WHO, 2022), dijital sağlık teknolojilerinin ruh sağlığı hizmetlerine entegrasyonunun öncelikli bir alan olduğunu belirtmektedir. Ancak bu entegrasyonun etik ve güvenlik standartlarına uygun biçimde gerçekleştirilmesi gerekmektedir.

\section{Alanyazındaki Boşluklar}

Yapılan sistematik literatür taraması sonucunda aşağıdaki boşluklar tespit edilmiştir:

\begin{enumerate}
  \item \textbf{Türkçe Kaynak Eksikliği:} HTP testinin Türkçe yorumlama kılavuzları sınırlıdır. Yavuzer (2009), Türkiye'de çocuk çizimleri üzerine kapsamlı çalışmalar yapmış olmakla birlikte, HTP'ye özgü sistematik bir yorumlama kılavuzu bulunmamaktadır. Mevcut kaynakların çoğu İngilizce literatürden çeviri niteliğindedir.

  \item \textbf{Dijital Araç Yokluğu:} Türkiye'de okul psikolojik danışmanlarına yönelik, çocuk çizimi analizi destekleyen dijital bir araç bulunmamaktadır. Uluslararası alanda ise sınırlı sayıda ticari yazılım mevcuttur; ancak bunların çoğu Türkçe desteği sunmamaktadır.

  \item \textbf{Yapay Zekâ Uygulamaları:} Uluslararası literatürde HTP analizi için yapay zekâ kullanımı araştırılmakta (Kim ve ark., 2024; Chen ve ark., 2023); ancak pratik uygulamalar henüz sınırlıdır. PICK (2024) çalışması, MLLM'lerin HTP analizinde umut verici sonuçlar ürettiğini raporlamıştır.

  \item \textbf{Kültürel Uyarlama:} HTP yorumlama kurallarının Türk kültürüne uyarlanması konusunda sistematik çalışmalar yetersizdir. Kültürlerarası psikoloji araştırmaları, çizim yorumlamalarının kültürel bağlama göre farklılık gösterebileceğini ortaya koymaktadır (Golomb, 2004; Thomas ve Silk, 1990).
\end{enumerate}

\section{Problemin Birey, Okul ve Toplum Açısından Önemi}

\subsection{Birey Düzeyinde}

Çocukların psikolojik sorunlarının erken dönemde tespit edilememesi, uzun vadeli olumsuz sonuçlara yol açabilmektedir. Costello, Mustillo, Erkanli, Keeler ve Angold (2003), tedavi edilmeyen çocukluk dönemi anksiyete ve depresyon belirtilerinin, ergenlik ve yetişkinlik döneminde daha ciddi psikolojik sorunlara dönüşebildiğini göstermiştir.

Kessler ve ark. (2005), yetişkinlik dönemindeki ruhsal bozuklukların yarısının 14 yaşından önce, \%75'inin ise 24 yaşından önce başladığını ortaya koymuştur. Bu bulgular, erken dönem tespitinin kritik önemini vurgulamaktadır.

\subsection{Okul Düzeyinde}

Psikolojik sorunlar, akademik başarıyı doğrudan etkilemektedir. McLeod ve Fettes (2007), çocukluk dönemi ruhsal sorunlarının okul başarısını \%20-30 oranında düşürdüğünü raporlamıştır. Erken müdahale programları, hem bireysel başarıyı artırmakta hem de sınıf içi davranış sorunlarını azaltmaktadır.

Durlak ve ark. (2011) tarafından yürütülen meta-analiz, okul temelli sosyal-duygusal öğrenme programlarının akademik başarıda 11 puanlık artış sağladığını göstermiştir.

\subsection{Toplum Düzeyinde}

Toplum ruh sağlığı perspektifinden bakıldığında, okul temelli önleme ve erken müdahale programları, maliyet-etkinlik açısından en verimli yaklaşımlar arasındadır. Knapp, McDaid ve Parsonage (2011), erken müdahale yatırımlarının uzun vadede 5-7 kat ekonomik geri dönüş sağladığını hesaplamıştır.

Dünya Sağlık Örgütü (WHO, 2021), küresel hastalık yükünün \%14'ünün ruhsal ve nörolojik bozukluklardan kaynaklandığını, bu yükün önemli bir kısmının erken müdahale ile önlenebileceğini belirtmektedir.

\newpage
